% GENERATED FILE. Edit canonical lessons, then run npm run generate:books.
% canonical-source-hash: fnv1a64:58309fc3b9045c28
% canonical-lessons: PA-C37-nahin-vaak, PA-C37-samajh-nahin, PA-C37-nahin-write, PA-C37-nahin-pata, PA-C37-nahin-checkpoint

\chapter{The Sentence That Says No}
\label{ch:pa-sentence-that-says-no}
\hlchaptermodality{37}

\begin{chapteropening}
\textbf{By the end of this chapter:} \emph{I can put nahin inside a sentence, contradict a statement, say I did not follow, and say I do not know.}

nahin has been in this book since chapter 1 and could only ever be a whole answer: both negative words were taught and no negated CLAUSE pattern was, so nothing in the corpus could be contradicted. The rule is one position -- nahin directly before the verb, nothing else moving -- and the two sentences built on it use the dative frame the track already owns from mainu ... pasand hai, so the repair language is new grammar for nobody.
\end{chapteropening}

\section[… nahī̃ …]{\pa{ਨਹੀਂ} inside a sentence}

\label{lesson:PA-C37-nahin-vaak}



\begin{quote}
From the first chapter of this book: the two words that answer a question. And from the third: the word for \textbf{I}.

\begin{itemize}
\raggedright
  \item \emph{Recall:} \emph{hā̃}, \emph{nahī̃} — then \emph{maiṁ}
\end{itemize}

You have been able to say \textbf{no} since page one. You have never been able to say that something \textbf{is not so}.
\end{quote}

\begin{cousinweb}
\textbf{\pa{ਨਹੀਂ}} goes \textbf{directly before the verb}, and nothing else in the sentence moves:

\begin{quote}
\pa{ਮੈਂ} \emph{(…)} \textbf{\pa{ਨਹੀਂ}} \emph{(verb)} — \textquotedblleft{}I do \textbf{not} \emph{(verb)}\textquotedblright{}
\end{quote}

That is the whole rule. Punjabi does not add a helper the way English adds \emph{do} — it drops one word in front of the verb and leaves the rest alone.

Because the position is fixed, every sentence you already own has a negative twin, and you can make it without learning anything else.
\end{cousinweb}

\subsection*{Your turn}
\begin{itemize}
\raggedright
  \item \emph{Say it:} \emph{nahī̃} alone, as the answer you already had.
  \item \emph{Build:} a sentence you can already say, then the same sentence with \emph{nahī̃} in front of its verb.
  \item \emph{Contrast:} the positive and the negative, back to back.
\end{itemize}

\begin{tcolorbox}[breakable,colback=teal!4,colframe=teal!35!black,title={Before you move on}]
Where does \textbf{\pa{ਨਹੀਂ}} stand in a sentence? (Directly before the verb.) What else has to change? (Nothing.)
\end{tcolorbox}
\section[mainū̃ samajh nahī̃ āī]{\pa{ਮੈਨੂੰ} \pa{ਸਮਝ} \pa{ਨਹੀਂ} \pa{ਆਈ} — I did not understand}

\label{lesson:PA-C37-samajh-nahin}



\begin{quote}
One lesson back, where \textbf{\pa{ਨਹੀਂ}} stands. From the verb chapter, the verb \textbf{to understand}. And two chapters back, the word for \textbf{to me}.

\begin{itemize}
\raggedright
  \item \emph{Recall:} \emph{nahī̃} before the verb, \emph{samajhṇā}, \emph{mainū̃}
\end{itemize}
\end{quote}

\begin{cousinweb}
\begin{itemize}
\raggedright
  \item \textbf{\pa{ਮੈਨੂੰ} \pa{ਸਮਝ} \pa{ਨਹੀਂ} \pa{ਆਈ}} — \emph{mainū̃ samajh nahī̃ āī} — \textquotedblleft{}understanding did not come to me\textquotedblright{}
\end{itemize}

Punjabi does not make you the subject of understanding. It makes understanding something that \textbf{arrives}, and you are the one it arrives at — which is exactly the shape you already met in \textbf{\pa{ਮੈਨੂੰ} … \pa{ਪਸੰਦ} \pa{ਹੈ}}, \textquotedblleft{}X is pleasing to me\textquotedblright{}.

So this is not a new grammar. It is the dative frame you own, with a different noun in it and \textbf{\pa{ਨਹੀਂ}} in the slot the last lesson gave you.
\end{cousinweb}

\begin{culture}
This is the sentence that buys you a second attempt. Without it a learner whose conversation has run away has only silence or an English apology, and the conversation ends there. It is worth more than any noun in this book, and it should arrive fast enough to be said without assembling it.
\end{culture}

\subsection*{Your turn}
\begin{itemize}
\raggedright
  \item \emph{Say it:} \emph{mainū̃ samajh nahī̃ āī}, slowly, then at speed.
  \item \emph{Contrast:} \emph{mainū̃ … pasand hai} and \emph{mainū̃ samajh nahī̃ āī} — one frame.
  \item \emph{Run through:} somebody says something too fast; reply with it.
\end{itemize}

\begin{tcolorbox}[breakable,colback=teal!4,colframe=teal!35!black,title={Before you move on}]
Say that you did not understand. Which earlier frame is it built on, and who is the subject of the sentence? (\textbf{\pa{ਮੈਨੂੰ} … \pa{ਪਸੰਦ} \pa{ਹੈ}}; the understanding is, not you.)
\end{tcolorbox}
\section[nahī̃]{\pa{ਨਹੀਂ} reaches the page}

\label{lesson:PA-C37-nahin-write}



\begin{quote}
Say where \textbf{\pa{ਨਹੀਂ}} stands in a sentence before anything is uncovered.
\end{quote}

\subsection*{Script — four pieces, none new}
Uncover \textbf{\pa{ਨਹੀਂ}} and name it in order:

\begin{quote}
\textbf{\pa{ਨ}} — \emph{na} · \textbf{\pa{ਹ}} — \emph{ha} · the \textbf{\pa{ੀ}} stroke · the \textbf{\pa{ਂ}} above
\end{quote}

The last of those is the bindi, and it is doing real work: it puts the vowel through the nose. A \textbf{\pa{ਨਹੀ}} without it is not this word.

This is the first of seven chapters, and it is worth saying now what holds for almost all of them: the words this book most needed were already writable. Every sign here was taught long ago.

\subsection*{Writing — guided copy, model in view}
Keep the model in view and copy \textbf{\pa{ਨਹੀਂ}} once. Check the bindi is there.

\begin{tcolorbox}[breakable,colback=teal!4,colframe=teal!35!black,title={Before you move on}]
Write \textbf{\pa{ਨਹੀਂ}} from memory, then uncover and compare. Which mark would change the word if you left it off? (The \textbf{\pa{ਂ}} above.)
\end{tcolorbox}
\section[mainū̃ nahī̃ patā]{\pa{ਮੈਨੂੰ} \pa{ਨਹੀਂ} \pa{ਪਤਾ} — I don't know}

\label{lesson:PA-C37-nahin-pata}



\begin{quote}
Two lessons back, the sentence that says you did not follow. And one lesson back, that word on the page.

\begin{itemize}
\raggedright
  \item \emph{Recall:} \emph{mainū̃ samajh nahī̃ āī}, and \textbf{\pa{ਨਹੀਂ}} written
\end{itemize}
\end{quote}

\begin{cousinweb}
\begin{quote}
\textbf{\pa{ਮੈਨੂੰ} \pa{ਨਹੀਂ} \pa{ਪਤਾ}} — \emph{mainū̃ nahī̃ patā} — \textquotedblleft{}I don't know\textquotedblright{}
\end{quote}

\textbf{\pa{ਪਤਾ}} is the new word, and it is a \textbf{noun} — the knowledge itself, not an act of knowing. So the sentence says, word for word, \emph{to me not knowledge}. The same dative frame again: the third sentence in three lessons to put \textbf{\pa{ਮੈਨੂੰ}} at the front and leave you out of the subject slot.

Punjabi keeps these two admissions apart, and so should you:

\begin{itemize}
\raggedright
  \item \textbf{\pa{ਮੈਨੂੰ} \pa{ਸਮਝ} \pa{ਨਹੀਂ} \pa{ਆਈ}} — I did not \textbf{follow} you
  \item \textbf{\pa{ਮੈਨੂੰ} \pa{ਨਹੀਂ} \pa{ਪਤਾ}} — I do not \textbf{have the fact}
\end{itemize}
\end{cousinweb}

\subsection*{Your turn}
\begin{itemize}
\raggedright
  \item \emph{Say it:} \emph{mainū̃ nahī̃ patā}, flat, as a whole reply.
  \item \emph{Contrast:} the two admissions, back to back.
  \item \emph{Run through:} asked a question you cannot answer, give the right one of the two.
\end{itemize}

\begin{tcolorbox}[breakable,colback=teal!4,colframe=teal!35!black,title={Before you move on}]
Say you do not know. Then say you did not follow. What kind of word is \textbf{\pa{ਪਤਾ}}? (A noun — the knowledge itself.)
\end{tcolorbox}
\section[Practice]{Saying no to a sentence, all together}

\label{lesson:PA-C37-nahin-checkpoint}



\begin{quote}
No new words. One run through what this chapter added.
\end{quote}

\begin{grammarlens}[title={What you've built}]
\begin{itemize}
\raggedright
  \item \textbf{\pa{ਨਹੀਂ}} directly before the verb, and nothing else moving.
  \item \textbf{\pa{ਮੈਨੂੰ} \pa{ਸਮਝ} \pa{ਨਹੀਂ} \pa{ਆਈ}} — I did not follow you.
  \item \textbf{\pa{ਮੈਨੂੰ} \pa{ਨਹੀਂ} \pa{ਪਤਾ}} — I do not have the fact.
  \item \textbf{\pa{ਨਹੀਂ}} on the page, bindi and all.
\end{itemize}

\textbf{\pa{ਨਹੀਂ}} has been in this book since its first chapter and could only ever be a whole answer. It can now be put inside a sentence, which means every statement you meet can be contradicted and every statement you make can be denied.
\end{grammarlens}

\subsection*{Your turn}
\begin{itemize}
\raggedright
  \item \emph{Say it:} someone states something you disagree with. Contradict it.
  \item \emph{Say it:} the two admissions, in the order you would really need them.
  \item \emph{Write it:} \textbf{\pa{ਨਹੀਂ}} from memory, then check the bindi.
  \item \emph{Run through:} refuse a yes-or-no question, then say why you cannot answer it.
\end{itemize}

\begin{tcolorbox}[breakable,colback=teal!4,colframe=teal!35!black,title={Before you move on}]
Contradict a statement, then admit you do not know. Next chapter: the word for \textbf{and}, which this book has never had either.
\end{tcolorbox}
